\section{Sub-agent 委派与并发控制}

\subsection{Sub-agent：用 step.invoke() 派生子 Agent}

有时 Agent 需要执行一个足够大的任务，大到会撑爆当前 context window——例如重构一个文件、调研一个课题、撰写一篇文档。对于像 OpenClaw 这样运行在单线程对话中的通用 Agent，长时间运行的会话可能因 context window 膨胀而出现问题。

解决方案：\textbf{派生 Sub-agent}。Utah 提供了 \texttt{delegate\_task} tool。当主 Agent 调用它时，它使用 \texttt{step.invoke()} 启动一个\textbf{完全独立的 Agent 函数运行实例}。

\begin{lstlisting}[caption={主 Agent 委派子任务},language=Java]
// In the main agent loop, when delegate_task is called:
const subResult = await step.invoke("sub-agent", {
  function: subAgent,
  data: {
    task: tc.arguments.task,
    subSessionKey: `sub-${sessionKey}-${Date.now()}`,
  },
});
\end{lstlisting}

Sub-agent 的关键特征：

\begin{itemize}
    \item 拥有自己独立的 context window（fork 父会话上下文到独立子会话）
    \item 使用相同的 Tool 集合（但\textbf{去除了 \texttt{delegate\_task}}——禁止递归派生）
    \item 拥有自己的重试策略和 step 级可观测性
    \item 完成后向父 Agent 返回一个摘要
    \item 父 Agent 只看到一个 Tool 结果："这是我完成的内容"
\end{itemize}

\begin{lstlisting}[caption={Sub-agent 函数定义（简化版）},language=Java]
export const subAgent = inngest.createFunction(
  { id: "agent-sub-agent", retries: 1 },
  { event: "agent.subagent.spawn" },
  async ({ event, step }) => {
    const { task, subSessionKey } = event.data;

    const agentLoop = createAgentLoop(task, subSessionKey, {
      tools: SUB_AGENT_TOOLS, // No delegate_task
      isSubAgent: true,
    });

    return await agentLoop(step);
  }
);
\end{lstlisting}

\begin{importantbox}{Sub-agent 的优雅之处}
\texttt{step.invoke()} 本身就是 Inngest 的原生能力——它用于在一个函数中调用另一个函数、等待结果、然后继续执行。Sub-agent 并不需要任何特殊的 ``Agent-to-Agent 协议''。它只是函数调用函数。编排由基础设施自动处理。
\end{importantbox}

\subsection{Singleton 并发：一次只处理一个对话}

每个通信渠道使用特定于渠道的 session key 来定义``对话''的边界。对于单线程渠道（如 Telegram），session key 是 chat ID；对于支持线程的平台（如 Slack），session key 包含 channel 和 thread 信息。

当多条消息在同一对话中发送时，你不希望第一个 Agent 循环继续运行，然后第二个循环再独立响应——你希望 Agent 拥有两条消息的上下文。Utah 选择的策略是：\textbf{取消当前运行，带着完整上下文重启}。

实现只需要一行配置：

\begin{lstlisting}[caption={Singleton 并发配置},language=Java]
singleton: { key: "event.data.sessionKey", mode: "cancel" },
\end{lstlisting}

这一行做了两件事：

\begin{enumerate}
    \item \textbf{Singleton 并发}：以 \texttt{sessionKey} 为键，同一时刻每个对话只有一个 Agent 运行实例。没有竞态条件，没有交错响应。
    \item \textbf{新消息时取消}：如果用户在 Agent 处理中发送新消息，当前运行被取消，新运行以最新消息启动。
\end{enumerate}

\begin{warningbox}{Singleton 并发的权衡}
取消+重启策略简洁，但有代价：任何正在进行中的工作都会丢失。新运行从持久化的 session 状态恢复，但这个过渡并不无缝。这是文章作者坦承的一个尚未完美解决的问题。在传统架构中，你需要自己构建每用户队列、管理锁和处理取消逻辑。Inngest 将其简化为一行配置，但底层问题仍然存在。
\end{warningbox}

\subsection{本章小结}

Sub-agent 委派利用 Inngest 原生的 \texttt{step.invoke()} 实现，无需额外的 Agent 间通信协议。Singleton 并发控制通过一行声明式配置即可解决在传统架构中需要大量代码的竞态问题。这两个机制共同展示了 harness 思维的核心优势：\textbf{将复杂的编排问题下沉到基础设施层}。


